\section{The finite Christoffel return to the original tensor Gamma source}
\label{sec:gcr-finite-christoffel-return}

This calculation retains the complete actual quartet, its order, the
original coordinate, and the same coefficient observation throughout:
\[
\begin{gathered}
 \rho=\tfrac12+\delta+i\gamma,\quad 0<\delta<\tfrac12,\quad
 \gamma>2,\quad m\geq1,\qquad
 h(s)=\prod_{\lambda\in\{\rho,\bar\rho,1-\rho,1-\bar\rho\}}
                  (s-\lambda)^m,\quad v_h=(2\xi)/h,\\
 k=4l+1,\quad l\geq0,\quad c=k/2,\quad
 e_k=1+k(m-1),\quad q=e_k(k+1)^2,\\
 \lambda_{a,b}=c+(2a-k)\delta+i(2b-k)\gamma,\qquad
 \chi(S)=\prod_{a,b=0}^k(S-\lambda_{a,b})^{e_k},
 \quad E_\chi=\mathbb C[S]/(\chi),\quad S=c+iy.
\end{gathered}\tag{GCR1}
\]
The symbol $e_k$ here is precisely the integer denoted $\ell_k$ in
TWA1; it is given a different letter only to keep it distinct from
the integer $l=(k-1)/4$. Both integers retain their displayed values.
Every divided derivative means $P^{[r]}=P^{(r)}/r!$.
Let $\mathcal P_N=\mathbb C[S]_{\leq N}$, and let
$J_{\chi,N}:\mathcal P_N\to E_\chi$ be literal polynomial remainder.
For $N\geq q-1$ it is onto and
\[
 \ker J_{\chi,N}=\chi\mathcal P_{N-q},\qquad
 \mathcal P_{-1}=\{0\}.
\tag{GCR2}
\]
The original full Taylor unit and the full original observation
$\Lambda:E_\chi\to\mathcal B=\operatorname{im}\Lambda$ are unchanged.
In particular the present source calculation does not substitute a
simple-root packet or remove any summand of $\ker\Lambda$.

\subsection{The exact Gamma factor, with its original mass}

Write
\[
\begin{gathered}
 d\sigma(y)=\frac{|\Gamma(1/4+iy/2)|^2}{2\pi}\,dy,\qquad
 c_a=2^{1-2a}\Gamma(2a),\\
 dm_{0,k}(y)=\frac{(2\pi)^{k/2}}{c_{k/4}}\,
                \frac{|\Gamma(k/4+iy/2)|^2}{2\pi}\,dy,\\
 t_j=2j+\tfrac12,\quad
 F_l(y)=\prod_{j=0}^{l-1}(y+it_j),\quad
 f_l(S)=\prod_{j=0}^{l-1}(S-c-t_j),\\
 \beta_k=\frac{(2\pi)^{k/2}}{4^l c_{k/4}}
         =\frac{(2\pi)^{k/2}}{\sqrt2\,\Gamma(k/2)}.
\end{gathered}\tag{GCR3}
\]
Empty products are one. Repeated application of
$\Gamma(z+1)=z\Gamma(z)$ gives
\[
 \left|\Gamma(\tfrac14+l+iy/2)\right|^2
 =4^{-l}\prod_{j=0}^{l-1}(y^2+t_j^2)
                      |\Gamma(\tfrac14+iy/2)|^2.
\]
Also $f_l(c+iy)=i^lF_l(y)$, with $|i^l|=1$.
It follows that the exact equality of measures is
\[
 \boxed{\quad dm_{0,k}(y)=\beta_k
                 |f_l(c+iy)|^2\,d\sigma(y).\quad}
\tag{GCR4}
\]
This preserves the masses
$\int d\sigma=\sqrt{2\pi}$ and
$\int dm_{0,k}=(2\pi)^{k/2}$. In particular $\beta_1=1$.
The equality is a change in the displayed density by a proved
polynomial factor; it does not divide either original measure
by its mass.

Let $T_{f,N}:\mathcal P_N\to\mathcal P_{N+l}$ be multiplication
by the monic polynomial $f=f_l$. If $H^\sigma_N$ is the Gram
in the increasing powers $1,S,\ldots,S^N$, then
\[
 H^f_N=T_{f,N}^*H^\sigma_{N+l}T_{f,N},\qquad
 H^0_N=\beta_k H^f_N,\qquad
 \|T_{f,N}P\|_\sigma^2=\beta_k^{-1}\|P\|_{0,k}^2.
\tag{GCR5}
\]
All these maps are injective on the stated polynomial spaces.
The coefficient change $P(S)\mapsto P(c+iy)$ is invertible with
inverse $Q(y)\mapsto Q((S-c)/i)$; its determinant is
$\prod_{j=0}^N i^j$, of absolute value one. It therefore preserves
each relative determinant in GCR5. No new coordinate for the
original quotient or observation is required.

\subsection{The product-jet map, including collisions}

We first prove the finite construction for any monic $f$ of degree
$l$, allowing common roots with the given monic $\chi$ of degree
$q$. Put $d=q+l$, and take the union of their distinct roots.
At a root $\zeta$ write
\[
 a_\zeta=\operatorname{ord}_\zeta f,\qquad
 b_\zeta=\operatorname{ord}_\zeta\chi,\qquad
 u_\zeta(S)=f(S)/(S-\zeta)^{a_\zeta}.
\]
Here $u_\zeta(\zeta)\ne0$, even when $a_\zeta b_\zeta>0$.
The product polynomial $f\chi$ has multiplicity
$a_\zeta+b_\zeta$ at that root. Define full Taylor maps
$\mathcal E_\chi$, $\mathcal E_f$, and $\mathcal E_{f\chi}$ on
their increasing monomial remainder bases, with all derivatives
of orders strictly below those multiplicities. The polynomial
division algorithm, or repeated factor divisibility, proves the
exact sequence
\[
 0\longrightarrow E_\chi
   \xrightarrow{\ I_f\ }E_{f\chi}
   \xrightarrow{\ r_f\ }E_f\longrightarrow0,\qquad
 I_f[z]=[fz],\quad r_f[P]=[P]\pmod f.
\tag{GCR6}
\]
Indeed $I_f[z]=0$ means $f\chi$ divides $fz$, and cancellation
in $\mathbb C[S]$ gives $\chi\mid z$. Moreover a remainder in
$\ker r_f$ is represented by $fQ$, and changing $Q$ by a multiple
of $\chi$ has exactly the required effect. These statements do
not use $\gcd(f,\chi)=1$.

At each $\zeta$, order the product jets into its first $a_\zeta$
entries, denoted $X$, and its next $b_\zeta$ entries, denoted $Y$.
Collecting the $X$ entries over all roots gives $l$ observations,
precisely the full jets of $f$. Collecting the $Y$ entries gives
$q$ observations, which are higher derivatives at shared roots.
Multiplication by $f$ is the following exact map in these coordinates:
\[
 \mathcal E_{f\chi}I_f\mathcal E_\chi^{-1}z
       =\binom{0}{U_fz},\qquad
 (U_fz)_{\zeta,r}
       =\sum_{s=0}^r u_\zeta^{[r-s]}(\zeta)z_{\zeta,s},
 \quad 0\leq r<b_\zeta.
\tag{GCR7}
\]
Thus $U_f$ is a direct sum of invertible lower triangular
multiplication matrices, and
\[
 \det U_f=\prod_{\zeta:b_\zeta>0}
                          u_\zeta(\zeta)^{b_\zeta}\ne0.
\tag{GCR8}
\]
To check GCR7, expand $f=(S-\zeta)^{a_\zeta}u_\zeta$ and $z$
in their Taylor series and take the coefficient of each power.
The first $a_\zeta$ coefficients vanish; the next $b_\zeta$
coefficients give exactly the displayed convolution. This
establishes the collision multiplicities, every nilpotent
coefficient and the actual map, without substituting a zero
resultant into an inverse.

For the particular $f_l$ and actual packet in GCR1, its roots
$\xi_j=c+t_j$ are distinct and real. Since $k=4l+1$ is odd,
$2b-k$ is an odd nonzero integer for every $b=0,\ldots,k$.
Consequently
\[
 |\operatorname{Im}(\lambda_{a,b}-\xi_j)|
       =|2b-k|\gamma\geq\gamma>2.
\tag{GCR9}
\]
There are no collisions on this entire stated parameter domain.
This is a direct proof from the original roots. In this case
$Y$ consists of the original full $\chi$ jets and
$U_f$ is literal multiplication by $f$ in those jets; GCR8 becomes
\[
 \det U_f=\prod_{a,b=0}^k f(\lambda_{a,b})^{e_k}
 =\prod_{a,b=0}^k\prod_{j=0}^{l-1}
       \big((2a-k)\delta-t_j+i(2b-k)\gamma\big)^{e_k}.
\tag{GCR10}
\]
Every factor is retained and nonzero. GCR6--GCR8 also specify
the exact construction if a different packet has shared roots.

\subsection{The joint kernel and its full conditional Gram}

Let $p_n(S)=i^n b_n((S-c)/i)$ be the monic polynomials for
$\sigma$ on the original line, with squared norms
\[
 \gamma_n=\sqrt{2\pi}\frac{(2n)!}{4^n},\qquad
 yb_n(y)=b_{n+1}(y)+n(n-\tfrac12)b_{n-1}(y).
\tag{GCR11}
\]
The absent $b_{-1}$ term is zero. For any finite ordered
list $\mathcal A$ of divided-jet observations, put
\[
 (K_{\mathcal A,L})_{\mu,\nu}
 =\sum_{n=0}^L
        \frac{\mathcal A_\mu(p_n)\,
                   \overline{\mathcal A_\nu(p_n)}}{\gamma_n}.
\tag{GCR12}
\]
Cross kernels use the first list on the first factor and the
second list on the second factor. These are full finite matrices.
If the list contains all jets of a degree-$r$ monic polynomial,
$L\geq r-1$ makes its observation onto: polynomials of degree
at most $r-1$ already interpolate every such jet. Its kernel
matrix is therefore positive definite.

Take the product jets from GCR7 at $L=N+l$, $N\geq q-1$,
and write their kernel as
\[
 K_{f\chi,N+l}
   =\begin{pmatrix}K_{XX}&K_{XY}\\K_{YX}&K_{YY}\end{pmatrix},
 \qquad
 C_N=K_{YY}-K_{YX}K_{XX}^{-1}K_{XY}>0.
\tag{GCR13}
\]
When $l=0$ the $X$ block is empty, its determinant is one,
and $C_N=K_{\chi,N}$. For $l>0$ positivity follows by completing
the square in the $X$ vector of the positive product kernel.
All cross terms in GCR13 are retained; in particular $C_N$
is not replaced by $K_{YY}$.

The minimum-norm source with a given finite observation $Az=z_0$
in a positive Gram $H$ is
\[
 R_{H,A}z_0=H^{-1}A^*(AH^{-1}A^*)^{-1}z_0.
\tag{GCR14}
\]
It has observation $z_0$. Its difference from any other vector
with that observation lies in $\ker A$, and is $H$-orthogonal
to the displayed vector. Expansion of the squared norm then
proves its minimality and its quotient Gram
$(AH^{-1}A^*)^{-1}$. This also proves uniqueness.

Apply GCR14 on $\mathcal P_{N+l}$ to the full product observation
with value $(0,U_f\mathcal E_\chi z)$. Every polynomial with
that observation is congruent to $fz$ modulo $f\chi$, so it is
divisible by $f$. Thus the entire constrained fibre is exactly
\[
 T_{f,N}\{P\in\mathcal P_N:J_{\chi,N}P=z\}.
\tag{GCR15}
\]
In particular the ambient minimum vector belongs to the range
of $T_{f,N}$; the preceding argument did not minimize over a
larger affine fibre. The bottom-right block of the inverse
matrix in GCR13 is $C_N^{-1}$, as direct block multiplication
shows. Consequently the full original quotient Grams satisfy
\[
 \boxed{\quad
 G_N^f=\mathcal E_\chi^*U_f^*C_N^{-1}U_f\mathcal E_\chi,
 \qquad G_N^0=\beta_k G_N^f.
 \quad}\tag{GCR16}
\]
The explicit source minimum vector is the inverse under $T_{f,N}$
of GCR14 for the product observation and value in GCR15.
Equivalently it is
$(H_N^f)^{-1}J_{\chi,N}^*G_N^f z$. These are equal because both
have that observation and are orthogonal to
$\chi\mathcal P_{N-q}$. Thus both the source injection and its
minimum-section inverse have been proved on their exact domains.

Taking determinants in GCR16 gives a finite Christoffel identity:
\[
 \det G_N^0
 =\beta_k^q|\det\mathcal E_\chi|^2|\det U_f|^2
                  \frac{\det K_{f,N+l}}{\det K_{f\chi,N+l}}.
\tag{GCR17}
\]
Here $K_{f,N+l}=K_{XX}$, including all multiplicities in the
collision construction. No determinant in this expression is
singular.

There is also an exact source and relation proof of this identity.
For any monic $g$ of degree $r$ and $L\geq r-1$, the increasing
frame
$1,\ldots,S^{r-1},g,gS,\ldots,gS^{L-r}$
has coefficient determinant one. Apply GCR14 to the full jets
of $g$ and take the block determinant in this frame. The result is
\[
 \det H^{|g|^2}_{L-r}
   =\left(\prod_{n=0}^L\gamma_n\right)
           \frac{\det K_{g,L}}{|\det\mathcal E_g|^2}.
\tag{GCR18}
\]
When $L=r-1$ its left side is the empty determinant one.
Apply the same frame argument to the original remainder map
$J_{\chi,N}$ with source weight $|f|^2$ to obtain
\[
 \det G_N^f
   =\frac{\det H^f_N}{\det H^{|f\chi|^2}_{N-q}}.
\tag{GCR19}
\]
The denominator is exactly the Gram of
$f\chi,f\chi S,\ldots,f\chi S^{N-q}$, pulled back to the
original relation coordinates. Equations GCR18--GCR19 give
GCR17. In the presence of collisions their constants agree
because
\[
 \frac{|\det\mathcal E_{f\chi}|^2}
      {|\det\mathcal E_f|^2|\det\mathcal E_\chi|^2}
       =\prod_{\zeta:b_\zeta>0}|u_\zeta(\zeta)|^{2b_\zeta}
       =|\det U_f|^2.
\tag{GCR20}
\]
For completeness, the divided-jet Vandermonde determinant is
$\prod_{\zeta<\eta}(\eta-\zeta)^{r_\zeta r_\eta}$
in increasing blocks of multiplicities $r_\zeta$.
It follows by merging distinct evaluation points within each
block after dividing their within-block Vandermonde factors;
Taylor expansion leaves precisely the divided derivatives.
Substitution of $r_\zeta=a_\zeta+b_\zeta$, $a_\zeta$, and
$b_\zeta$ cancels the same-type exponents and leaves
$a_\zeta b_\eta+b_\zeta a_\eta$ for each distinct pair.
Expanding $u_\zeta(\zeta)$ gives those same absolute-square
factors, proving GCR20 with no unaccounted factorial.

\subsection{The exact signed return is a determinant of size $2l$}

Specialize again to GCR1--GCR4, where GCR9 proves coprimality.
Put $K_N=K_{\chi,N}$ and let
\[
\begin{split}
 (U_N)_{\mu,j}
   &=\frac{\mathcal E_{\chi,\mu}(p_{N+j})}
                                  {\sqrt{\gamma_{N+j}}},
                    \qquad 1\leq j\leq l,\\
 V_N&=K_{\chi,f;N+l}\,K_{f,N+l}^{-1/2},\qquad
 W_N=(\,U_N\ \ V_N\,),\qquad
 J_l=\operatorname{diag}(I_l,-I_l).
\end{split}\tag{GCR21}
\]
The $q$ observations $\mathcal E_{\chi,\mu}$ here act on arbitrary
polynomials by their full divided jets. $U_N$ is a matrix of
new orthonormal source columns; it is distinct from the
multiplication matrix $U_f$ in GCR7. Explicitly the cross kernel
in this formula is the $q$ by $l$ matrix
\[
 (K_{\chi,f;L})_{\mu,j}
    =\sum_{n=0}^L
       \frac{\mathcal E_{\chi,\mu}(p_n)\,
                         \overline{p_n(\xi_j)}}{\gamma_n},
       \qquad 0\leq j<l.
\]
Its notation with a comma distinguishes it from the square
full product-jet kernel $K_{f\chi,L}$. Directly subtracting the
finite sums in GCR12 and then using GCR13 proves
\[
 C_N=K_N+U_NU_N^*-V_NV_N^*
        =K_N+W_NJ_lW_N^*.
\tag{GCR22}
\]
The determinant identity $\det(I+AB)=\det(I+BA)$ follows, for
rectangular $A,B$, by computing the determinant of
$\left(\begin{smallmatrix}I&A\\-B&I\end{smallmatrix}\right)$
using either diagonal block. Therefore
\[
 d_N:=\frac{\det C_N}{\det K_N}
       =\det\!\left(I_{2l}+J_lW_N^*K_N^{-1}W_N\right)>0.
\tag{GCR23}
\]
For $l=0$ this is one. The small matrix in GCR23 need not itself
be Hermitian; its determinant is strictly positive by the
first equality. A positive factorization retaining its mixed
entries is
\[
\begin{gathered}
 A_N=I_l+U_N^*K_N^{-1}U_N>0,\\
 B_N=I_l-V_N^*(K_N+U_NU_N^*)^{-1}V_N>0,\\
 d_N=\det A_N\det B_N,\\
 B_N=I_l-V_N^*K_N^{-1}V_N
       +V_N^*K_N^{-1}U_N A_N^{-1}U_N^*K_N^{-1}V_N .
\end{gathered}\tag{GCR24}
\]
The last identity is checked by multiplying the proposed inverse
$K_N^{-1}-K_N^{-1}U_N A_N^{-1}U_N^*K_N^{-1}$
by $K_N+U_NU_N^*$. Positivity of $B_N$ follows from the
Schur complement of the positive joint kernel in GCR13,
or equivalently from $C_N>0$ after congruence by
$(K_N+U_NU_N^*)^{-1/2}$. Its mixed product cannot be discarded.

Let $(N_0,N_1,N_2,N_3)=(q-1,q,2q-1,2q)$ and
$(\epsilon_0,\epsilon_1,\epsilon_2,\epsilon_3)=(1,1,-1,-1)$.
Define the two original signed quotient expressions
\[
 F_{0,k}=\sum_{a=0}^3\epsilon_a\log\det G_{N_a}^0,\qquad
 F_{\sigma,k}=\sum_{a=0}^3\epsilon_a\log\det G_{N_a}^\sigma.
\]
The fixed-source formula is
$G_N^\sigma=\mathcal E_\chi^*K_N^{-1}\mathcal E_\chi$.
Subtract it from GCR16 before summing to obtain
\[
 \log\det G_N^0-\log\det G_N^\sigma
      =q\log\beta_k+2\log|\det U_f|-\log d_N.
\tag{GCR25}
\]
Only after all four signs are retained do the first two terms
cancel. Thus the required return is the single explicit signed
object
\[
 \boxed{\quad
 \mathfrak T_k:=F_{0,k}-F_{\sigma,k}
    =-\sum_{a=0}^3\epsilon_a\log
       \det\!\left(I_{2l}
           +J_lW_{N_a}^*K_{N_a}^{-1}W_{N_a}\right).
 \quad}\tag{GCR26}
\]
Its complete mixed part is displayed in GCR24.
Equivalently it is the logarithm of the product of the two
positive endpoint determinants in the denominator-sign pair
divided by the product in the numerator-sign pair.
The four matrices use the original four source degrees and
the same packet; the construction introduces no independent
large-degree limit.

Using the proved finite identities TWA23--TWA24 gives
\[
 \mathcal R_k^0-\mathcal R_k^\sigma=-\mathfrak T_k
       =\sum_{a=0}^3\epsilon_a\log d_{N_a},\qquad
 \Delta_k^\Gamma=\delta_k^\sigma-\mathfrak T_k.
\tag{GCR27}
\]
Thus the determinant mass-rate result
$\delta_k^\sigma=o(kq)$ removes exactly that original-source
term at this scale. The remaining term is the four explicit
matrices in GCR26, including their mixed products. Formula
GCR27 does not assert that this term tends to zero.

\subsection{A finite determinant bound with proved moving constants}

The following bound uses the original full source dimension at
each endpoint and propagates its whole logarithmic deficit.
For $l\geq1$ and any $N\geq0$ set $n=N+1$ and
\[
\begin{gathered}
 \mathscr U_{l,N}
    =\prod_{j=0}^{l-1}\bigl(4(N+j+1)^2+t_j^2\bigr),\\
 D_{l,N}
    =\prod_{a=0}^{N}\prod_{r=0}^{2l-1}(a+r+\tfrac12),\\
 X_{l,N}=n\log\mathscr U_{l,N}-\log D_{l,N}.
\end{gathered}\tag{GCR28}
\]
In the orthonormal form of GCR11, multiplication by $y$ has
off-diagonal coefficients $\sqrt{a(a-1/2)}\leq a$.
On a polynomial of degree $d$ its output has degree at most
$d+1$. The corresponding symmetric truncated Jacobi matrix
has every absolute row sum at most $2(d+1)$, so its operator
norm is at most $2(d+1)$. To verify this last elementary bound,
for a real symmetric matrix use
$2|x_i x_j|\leq |x_i|^2+|x_j|^2$ in its quadratic form;
the same proof applies to complex vectors using absolute values.
Consequently
\[
 \|yP\|_\sigma^2\leq4(d+1)^2\|P\|_\sigma^2,\qquad
 \|(y+it)P\|_\sigma^2
      =\|yP\|_\sigma^2+t^2\|P\|_\sigma^2.
\]
The second identity has no cross term because
$\langle P,yP\rangle_\sigma$ is real. Apply it successively
at degrees $N,N+1,\ldots,N+l-1$ and $t=t_0,\ldots,t_{l-1}$.
The exact coordinate map in GCR5 yields
\[
 0<H_N^f\preceq\mathscr U_{l,N}H_N^\sigma .
\tag{GCR29}
\]

The original tensor Gamma monic norms, with their original mass,
are
$\gamma_{k,a}=(2\pi)^{k/2}a!\Gamma(a+k/2)/\Gamma(k/2)$.
Dividing by the factor $\beta_k$ in GCR3 and then by
$\gamma_a=\sqrt2\,a!\Gamma(a+1/2)$ gives
\[
 \frac{\gamma_{k,a}}{\beta_k\gamma_a}
     =\frac{\Gamma(a+2l+1/2)}{\Gamma(a+1/2)}
     =\prod_{r=0}^{2l-1}(a+r+\tfrac12).
\]
Each change from monomials to its monic orthogonal frame has
determinant one. It follows that
\[
 \frac{\det H_N^f}{\det H_N^\sigma}=D_{l,N},\qquad
 X_{l,N}\geq0.
\tag{GCR30}
\]
Thus the factorial product and the source bound have the same
source degree $N$ and the same source dimension $n=N+1$;
the ambient degree $N+l$ in GCR5 has not been used as the
dimension in this determinant ratio.

Here is the full propagation to any onto rank-$q$ observation
$J$ on this source. Put
$E=(H_N^\sigma)^{-1/2}H_N^f(H_N^\sigma)^{-1/2}/
\mathscr U_{l,N}$, so $0<E\preceq I_n$ and
$\det E=\exp(-X_{l,N})$.
Choose an orthonormal basis of
$(H_N^\sigma)^{1/2}\ker J$ and its orthogonal complement.
An explicit orthonormal frame of the latter is
\[
 Q=(H_N^\sigma)^{-1/2}J^*
        \bigl(J(H_N^\sigma)^{-1}J^*\bigr)^{-1/2}.
\tag{GCR31}
\]
Indeed $Q^*Q=I_q$, and it is orthogonal to every vector in the
stated kernel. In this unitary source frame write
$E=\left(\begin{smallmatrix}E_{KK}&E_{KC}\\
E_{CK}&E_{CC}\end{smallmatrix}\right)$.
The minimum quotient in these complementary coordinates is
the exact Schur complement
$E_{CC}-E_{CK}E_{KK}^{-1}E_{KC}$. It is congruent to
\[
 (G_N^\sigma)^{-1/2}
       (G_N^f/\mathscr U_{l,N})(G_N^\sigma)^{-1/2},
\]
where $G_N^\sigma=(J(H_N^\sigma)^{-1}J^*)^{-1}$.
This follows either from GCR14 or by checking that the source
map $(H_N^\sigma)^{-1/2}Q$ has observation
$(G_N^\sigma)^{-1/2}$. The Schur complement is at most $I_q$.
Furthermore $\det E_{KK}\leq1$, so its determinant is at least
$\det E$ by the block determinant formula. At zero kernel
dimension its empty determinant is one and the same statement
holds. This proves, for the original remainder observation,
\[
 -X_{l,N}\leq
 \log\det G_N^f-\log\det G_N^\sigma
                      -q\log\mathscr U_{l,N}\leq0.
\tag{GCR32}
\]
All relation directions are present in $E_{KK}$; none has been
completed away or omitted from the determinant.

Put
\[
 Z_k=q\sum_{a=0}^3\epsilon_a\log\mathscr U_{l,N_a}.
\]
GCR32 at each of the four source degrees, together with the
cancellation of $\beta_k^q$, proves the finite signed interval
\[
 \boxed{\quad
 Z_k-X_{l,q-1}-X_{l,q}
      \ \leq\ \mathfrak T_k\
      \leq Z_k+X_{l,2q-1}+X_{l,2q}.
 \quad}\tag{GCR33}
\]
Every quantity on its two sides is the explicit finite product
in GCR28; it does not involve an unknown condition number.

We calculate the moving limits of these finite products.
Let $d=2l$. Directly splitting each factor in $\mathscr U_{l,N}$
gives
\[
 \log\mathscr U_{l,N}=d\log(2n)+\eta_{l,n},\qquad
 0\leq\eta_{l,n}\leq l^2/n+l^3/n^2.
\tag{GCR34}
\]
For the bound use
$2\log(1+j/n)\leq2j/n$ and
$\log(1+t_j^2/(4(n+j)^2))\leq t_j^2/(4n^2)$;
their sums are no greater than the displayed two terms because
$t_j<2l$.
The concavity of $\log x$ on each interval $[a,a+1]$ gives
\[
 \sum_{a=0}^{n-1}\log(a+\tfrac12)
       \geq\int_0^n\log x\,dx=n\log n-n.
\]
For example its midpoint inequality follows by pairing the
points at equal distance from $a+1/2$ and applying concavity;
the improper integral at zero is finite. Its upper bound is
$\log n!\leq n\log n-n+1+\log n$, by comparison of the
increasing function $\log x$ with its adjacent integral sums.
For an integer $r\geq0$,
\[
 0\leq\sum_{a=0}^{n-1}
          \log\frac{a+r+1/2}{a+1/2}
   \leq r\sum_{a=0}^{n-1}(a+1/2)^{-1}
   \leq r\bigl(2+\log(2n)\bigr).
\]
For the last inequality keep the $a=0$ term $2$ and bound the
remaining decreasing sum by
$\int_0^{n-1}(x+1/2)^{-1}dx\leq\log(2n)$.
Summing in $0\leq r<d$ proves
\[
\begin{split}
 \log D_{l,N}&=dn(\log n-1)+\theta_{l,n},\\
 0\leq\theta_{l,n}
 &\leq d(1+\log n)
       +\tfrac{d(d-1)}2\bigl(2+\log(2n)\bigr).
\end{split}\tag{GCR35}
\]
Consequently
\[
 X_{l,N}=dn(1+\log2)+n\eta_{l,n}-\theta_{l,n}.
\tag{GCR36}
\]
This is an equality with the explicit bounded remainders in
GCR34--GCR35, not a fixed-degree kernel asymptotic.

On the four original degrees, $n\geq q\geq(k+1)^2$,
$q=O_m(k^3)$, and $l=(k-1)/4$. Dividing the upper bounds
in GCR34--GCR35 by $d$ and $dn$, respectively, proves
\[
 \frac{X_{l,N_a}}{2l(N_a+1)}\longrightarrow1+\log2.
\tag{GCR37}
\]
Also
\[
 \frac{Z_k}{kq}
 =\frac{2l}{k}\log
       \frac{q(q+1)}{(2q)(2q+1)}
      +\frac1k\sum_{a=0}^3\epsilon_a\eta_{l,N_a+1}
       \longrightarrow-\log2.
\tag{GCR38}
\]
The remainder tends to zero by its displayed bound and
$q\geq(k+1)^2$. Since
$(X_{l,q-1}+X_{l,q})/(kq)\to1+\log2$ and
$(X_{l,2q-1}+X_{l,2q})/(kq)\to2(1+\log2)$,
GCR33 proves the moving conclusion
\[
 \boxed{\quad
 -1-2\log2
 \leq\liminf_{\substack{k\to\infty\\k\equiv1\ (4)}}
              \frac{\mathfrak T_k}{kq}
 \leq\limsup_{\substack{k\to\infty\\k\equiv1\ (4)}}
              \frac{\mathfrak T_k}{kq}
 \leq2+\log2.
 \quad}\tag{GCR39}
\]
These finite source calculations establish
$\mathfrak T_k=O_m(kq)$. They neither establish a nonzero
limit nor imply $o(kq)$. With the already proved determinant
mass-rate return in GCR27 they also give
\[
 -2-\log2
 \leq\liminf_{\substack{k\to\infty\\k\equiv1\ (4)}}
            \frac{\Delta_k^\Gamma}{kq}
 \leq\limsup_{\substack{k\to\infty\\k\equiv1\ (4)}}
            \frac{\Delta_k^\Gamma}{kq}
 \leq1+2\log2.
\tag{GCR40}
\]
The exact signed matrices, rather than the endpoints of these
enclosures, remain the value of the original correction.

\subsection{The original unitful observation and both retained kernels}

For clarity the exact return to the original observation can be
made within either of the source forms $H_N^f,H_N^0$.
Let $R_N=(H_N^f)^{-1}J_{\chi,N}^*G_N^f$ be the minimum
section already identified in GCR16. Then
$J_{\chi,N}R_N=I_{E_\chi}$, and
\[
 \ker(\Lambda J_{\chi,N})
       =\chi\mathcal P_{N-q}\ \mathbin{\perp_{H_N^f}}\
                                  R_N(\ker\Lambda).
\tag{GCR41}
\]
Every vector in the left side has the unique decomposition
$P=(P-R_NJ_{\chi,N}P)+R_NJ_{\chi,N}P$.
The first term is in the first displayed summand by GCR2;
the second is in the second summand. Their orthogonality is
the minimum-section identity GCR14. Conversely both summands
are killed by $\Lambda J_{\chi,N}$, proving equality.
Multiplication by $f$ sends the first summand exactly to
$f\chi\mathcal P_{N-q}$ and the second to $fR_N(\ker\Lambda)$
inside $\mathcal P_{N+l}$ with its original $\sigma$ metric.
Equation GCR5 proves their orthogonality there as well.

If $I_K:\mathcal K\to E_\chi$ is the original inclusion
of $\mathcal K=\ker\Lambda$, the original kernel and boundary
forms for a quotient Gram $G$ are exactly
\[
 H_K=I_K^*GI_K,\qquad
 Q_{\mathcal B}=(\Lambda G^{-1}\Lambda^*)^{-1},\qquad
 L_{\mathcal B}=G^{-1}\Lambda^*Q_{\mathcal B}.
\tag{GCR42}
\]
Here $\Lambda$ has codomain its actual image, so it is onto;
$L_{\mathcal B}$ is a section, and $G$-orthogonality to
$\mathcal K$ follows from the same calculation as GCR14.
For any fixed algebraic lift $L_*$ of $\Lambda$, the square
frame $T_*=(I_K,L_*)$ is invertible. Its determinant identity is
\[
 |\det T_*|^2\det G
       =\det H_K\,\det Q_{\mathcal B}.
\tag{GCR43}
\]
To prove it, replace $L_*$ by $L_{\mathcal B}$.
Their difference has image in $\mathcal K$, so the frame
changes by a block triangular matrix with identity diagonal.
The resulting two blocks are orthogonal and have precisely
the Grams in GCR42. At either rank-zero endpoint the empty
determinant is one. This proves GCR43 in every boundary rank.
Using the same fixed $T_*$ at the four degrees cancels its
constant with the same four signs. Thus $F_{0,k}$ is literally
the complete $F_K(0)+F_B(0)$ in TWA24.

No multiplication by $U_f$ in GCR7 is an unproved identification
of the original full Taylor unit with the Gamma factor.
The original observation still acts on $z\in E_\chi$ as
$\Lambda z$. In the ambient constrained product jets it acts
by the proved map
\[
 (0,Y)\longmapsto
       \Lambda\,\mathcal E_\chi^{-1}U_f^{-1}Y.
\tag{GCR44}
\]
GCR7 shows that this composite sends the jet of $fP$ exactly
to $\Lambda J_{\chi,N}P$, with all original nilpotents and
unit coefficients retained. This establishes the actual
commuting map through the Gamma source, including both
relation kernels in GCR41 and the complete mixed conditional
Gram in GCR13. The calculation closes this finite comparison;
it makes no assertion that the remaining signed term in
GCR26 contradicts or supplies an actual zero of $\xi$.
